\documentclass[12pt]{article}
\usepackage{amsmath}
\usepackage{amsfonts}
\usepackage{amssymb}
\usepackage{geometry}
\usepackage{graphicx}
\geometry{a4paper, margin=0.8in}
\pagenumbering{gobble}

\title{Homework assignment for 02YMECH \\ Chapter: Central-force motion}
\begin{document}
\date{}
\maketitle
\vspace{-0.8in}

\begin{enumerate}
\item A point mass moves in the plane with polar coordinates $(r,\theta)$ relative to a fixed origin. Let
\[
r(t)=r_0+v_r\,t,\qquad \theta(t)=\Omega\,t+\tfrac{1}{2}\alpha t^2,
\]
where \(r_0>0\), \(v_r\), \(\Omega\), and \(\alpha\) are constants. 

\begin{enumerate}
  \item[(a)] Write the general expressions for the velocity and acceleration in polar coordinates.
  \item[(b)] Find the radial and angular components of the velocity and acceleration as functions of time.
  \item[(c)] Compute the speed \(v(t)=\|\mathbf v\|\) and its time derivative \(\dot v(t)\). State the condition under which \(\dot v\) is positive, negative, or zero.

  \item[(d)] Find the time(s) \(t^\ast\) (if any) at which the radial component of acceleration \(a_r\) vanishes. State the condition on parameters for such a time to exist (with \(t^\ast\ge 0\)).
\end{enumerate}
\item     A particle of mass $m$ moves under a central force $f(r)$.  
    Its trajectory in polar coordinates is observed to be
    \[
        r(\theta)=\frac{a}{1+b\cos(\theta)},
    \]
    where $a>0$ and $0<b<1$. 

    \begin{enumerate}
        \item Determine the central force $f(r)$ as a function of $r$.
        \item State whether this force is attractive or repulsive for $0<r<a$.
    \end{enumerate} 

    \item (Hard) Perform an explicit calculation of time average of potential energy over a one complete period for a particle moving in an elliptical orbit under a central inverse square law force field. Express the result in terms of the force constant of the field and the semi-major axis of the ellipse. Perform a similar calculation for the time average of kinetic energy. Compare your results and verify that they are consistent with the virial theorem.
    \item Two particles are moving under the influence of their mutual gravitational force describe circular orbits about one another with a period $\tau$. If they suddenly stopped in their orbits and allowed to gravitate toward each other, show that they would collide after a time interval $\tau/4\sqrt{2}$.


\end{enumerate}
\end{document}
